\section{Congruential Euler packets and their arithmetic}
\label{sec:arithmetic}

Two related constructions run through this section.  The congruential
Euler family is indexed by a modulus $N$ and an address $j$.  Its
exceptional dyadic member has $N=4$ and four addresses.  The classical
signed Euler numbers arise from $N=2$, $j=0$.  Keeping these indices
visible distinguishes the four-packet period theorems from the classical
edge truncations.

\subsection{Original coefficients, normalized packets, and precision}

\begin{definition}[Congruential Euler coefficients]\label{def:arith-euler}
For $N\ge1$ and $0\le j<N$, define rational coefficients by the formal identity
\begin{equation}\label{eq:arith-euler-egf}
 \sum_{d\ge0}E_d^{(N,j)}\frac{t^d}{d!}
 =\left(\sum_{n\ge0}\frac{t^{Nn}}{(Nn+j)!}\right)^{-1}.
\end{equation}
The reciprocal exists because its denominator has constant coefficient
$1/j!\ne0$.  Set
\[
 U_{N,j}(n)=E_{Nn}^{(N,j)},\qquad C_{N,j}(n)=\frac{U_{N,j}(n)}{j!}.
\]
Thus $C_{N,j}(0)=1$ and $U_{N,j}(0)=j!$.  All coefficients outside the
multiples of $N$ vanish.
\end{definition}

The normalized source and its inverse are
\begin{equation}\label{eq:arith-kernel}
 K_{N,j}(t)=j!\sum_{n\ge0}\frac{t^{Nn}}{(Nn+j)!},\qquad
 \mathcal P_{N,j}(t)=\sum_{n\ge0}C_{N,j}(n)\frac{t^{Nn}}{(Nn)!}.
\end{equation}
Their product is one.  Clearing factorials in its coefficient of $t^{Nn}$
gives the triangular system
\begin{equation}\label{eq:arith-triangular}
 \sum_{q=0}^{n}\binom{Nn+j}{Nq}C_{N,j}(q)=\delta_{n0}.
\end{equation}
In particular, for $n>0$,
\[
 C_{N,j}(n)=-\binom{Nn+j}{Nn}^{-1}
       \sum_{q=0}^{n-1}\binom{Nn+j}{Nq}C_{N,j}(q).
\]
The diagonal binomial coefficient is positive, so this recurrence determines
the sequence uniquely.  The Lean construction begins with this recurrence
and derives the reciprocal identity from it.  The resulting sequence is
then identified with the degree-indexed coefficients of
\eqref{eq:arith-euler-egf}, so both descriptions refer to the same family.

The relevant declarations are
\lean{KomorebiPaperIX.Congruential.triangular_inverse},
\lean{KomorebiPaperIX.Congruential.packet_is_actual_inverse}, and
\lean{KomorebiPaperIX.NishibuchiCompletion.generalizedEuler_at_packet}.
The cleared identity with the original positive-address exponential section
is \lean{KomorebiPaperIX.NishibuchiCompletion.original_generating_identity}.

\begin{proposition}[A denominator bound for every address]
\label{prop:arith-denominator}
For $N\ge1$, $0\le j<N$, and $n\in\N$, set
\[
 D_{N,j}(n)=\prod_{r=1}^{n}\binom{Nr+j}{Nr},\qquad D_{N,j}(0)=1.
\]
Then $D_{N,j}(n)C_{N,j}(n)\in\Z$.
\end{proposition}

\begin{proof}
Proceed by induction in the recurrence \eqref{eq:arith-triangular}.
For every $q<n$, the product $D_{N,j}(q)$ divides $D_{N,j}(n-1)$.
Multiplying the recurrence by $D_{N,j}(n-1)$ therefore makes every
previous term integral.  The diagonal factor completes the product
$D_{N,j}(n)$, proving the claim.  At $n=0$, both factors are one.
This gives a uniform denominator bound without requiring integrality of
the normalized coefficients themselves.  The Lean statements are
\lean{KomorebiPaperIX.Congruential.clearing_product} and
\lean{KomorebiPaperIX.Congruential.universal_denominator_clearing}.
\end{proof}

\begin{example}[The Bernoulli boundary address]\label{ex:arith-bernoulli}
At $N=2$, the address $j=1$ gives a second classical family.  Define the
Bernoulli numbers by
\[
 B(z)=\frac{z}{e^z-1}=\sum_{r\ge0}B_r\frac{z^r}{r!},
 \qquad B_0=1,\quad B_1=-\frac12.
\]
All quotients here are regular formal series.  The identity
\[
 2B(z)-B(2z)=\frac{z}{\sinh z}
            =\sum_{n\ge0}C_{2,1}(n)\frac{z^{2n}}{(2n)!}
\]
follows by putting the two fractions on a common denominator and using
\eqref{eq:arith-euler-egf}.  Comparing even coefficients yields
\[
 C_{2,1}(n)=(2-2^{2n})B_{2n}\qquad(n\in\N).
\]
This classical derivation restores the boundary example from the original
manuscript.  No dedicated Lean bridge for this Bernoulli specialization
was identified in the preserved sources.
\end{example}

\begin{definition}[Differences and valuation conventions]\label{def:arith-precision}
For a sequence $f$ and $h,k,n\in\N$, put
\[
 \Delta_h f(n)=f(n+h)-f(n),\qquad
 \Delta_h^0 f(n)=f(n),\qquad
 \Delta_h^{k+1}f(n)=\Delta_h^k f(n+h)-\Delta_h^k f(n).
\]
We use the extended valuation $v_p(0)=+\infty$ in mathematical statements.
Thus $v_p(x)\ge L$, including negative integers $L$, means
$x=0$ or the ordinary valuation of the nonzero rational $x$ is at least $L$.
For integral coefficients and $L\ge0$, it is equivalent to
$x\in p^L\Z_p$.  In expressions involving a step multiplier we write
$\nu_p(u)=v_p(u)$ for $u>0$ and set $\nu_p(0)=0$.  This finite convention
makes the displayed exponents meaningful at step zero, where all
positive-order differences vanish.
\end{definition}

The valuation convention at zero is encoded by
\lean{KomorebiPaperIX.GlobalEulerBoundary.precision}.  For $p$-adic
congruences, divisibility is interpreted in $\Z_p$; statements over other
coefficient rings specify the ring explicitly.  Field divisibility in
$\Q_p$ alone does not encode precision.

\subsection{The standard address window and its common reduction}

The addresses below the prime have unit diagonal coefficients.  In this
window the normalized packets are integral, and their reductions agree
at every row.

\begin{proposition}[Lucas reduction and address independence]
\label{prop:arith-lucas-window}
Let $p$ be prime, $m>0$, and $N=mp$.  For $0\le j<p$ and $n,q\in\N$,
\[
 \binom{mpn+j}{mpq}\equiv\binom{mn}{mq}\pmod p.
\]
Every $C_{mp,j}(n)$ belongs to $\Z_p$.  Its reduction is the unique
sequence $\overline C_m:\N\to\F_p$ satisfying
\[
 \sum_{q=0}^{n}\binom{mn}{mq}\overline C_m(q)=\delta_{n0}
 \quad\text{in }\F_p.
\]
Consequently, for any two addresses $0\le j,k<p$,
\[
 C_{mp,j}(n)\equiv C_{mp,k}(n)\pmod p\qquad(n\in\N).
\]
\end{proposition}

\begin{proof}
Lucas reduction separates the lowest base-$p$ digits, which are $j$ and
zero in the upper and lower arguments.  Their binomial contribution is
one, leaving the displayed reduced row.  On the diagonal the reduced
coefficient is one.  The recurrence \eqref{eq:arith-triangular} therefore
constructs an integral packet by successive division by units in $\Z_p$.
After reduction, every address satisfies the same triangular system over
$\F_p$.  Its diagonal entries are one, so its solution is unique.  This
proves both the common reduction and the congruence between addresses.
The row reduction is
\lean{KomorebiPaperIX.StandardWindow.lucas_row_reduction}.
The integral identification and the two reduction statements are
\lean{KomorebiPaperIX.WindowIndependence.standard_window_integral},
\lean{KomorebiPaperIX.WindowIndependence.coeffPadic_coe},
\lean{KomorebiPaperIX.WindowIndependence.residue_independence}, and
\lean{KomorebiPaperIX.WindowIndependence.two_window_addresses_congruent}.
\end{proof}

\subsection{The full Nishibuchi statement and its starting bound}

\begin{theorem}[Nishibuchi congruence]\label{thm:arith-nishibuchi}
Let $p$ be prime, let $m>0$, and suppose that $m\mid p-1$ or $m=2$.
Put $N=mp$ and $q=\operatorname{lcm}(2,p-1)$.  For every $0\le j<N$,
all $U_{N,j}(n)$ belong to $\Z_p$.  For $r\ge1$ and $n\ge r$,
\begin{equation}\label{eq:arith-nishibuchi}
 E_{Nn+qp^r}^{(N,j)}\equiv E_{Nn}^{(N,j)}\pmod{p^r}.
\end{equation}
Thus the positive sufficient starting index $n_0=r$ works uniformly.
If $m\mid p-1$, the stronger sufficient condition $r\le mn$ is valid.
\end{theorem}

This theorem covers the full parameter range and displacement of Conjecture
3.9 in the fixed version of Nishibuchi's paper used for the statement
comparison.  The corresponding Lean theorem is
\lean{KomorebiPaperIX.NishibuchiCompletion.conjecture_3_9}.
The stronger divisor-branch threshold is
\lean{KomorebiPaperIX.NishibuchiCompletion.divisor_period}.
The starting index is sufficient uniformly, although it need not be the
earliest one for a given tuple of parameters.

We first prove the difference estimates on the divisor branch, then derive
the period from a finite Newton expansion.

\begin{theorem}[All-address differences on the divisor branch]
\label{thm:arith-divisor-base}
Let $p$ be prime, $m>0$, $m\mid p-1$, $N=mp$, and $0\le j<N$.
Write $h=(p-1)/m$.  For every $u,k,n\in\N$,
\begin{align}
 U_{N,j}(n)&\in\Z_p,\label{eq:arith-divisor-integral}\\
 \Delta_{uh}^{k}U_{N,j}(n)&\in p^{\min(k,mn)}\Z_p,\label{eq:arith-divisor-bound}\\
 v_p\bigl(\Delta_{uh}^{k}C_{N,j}(n)\bigr)
   &\ge\min(k,mn)-v_p(j!).\label{eq:arith-divisor-normalized}
\end{align}
The factorial correction is part of the normalized statement, including
addresses $j\ge p$.
\end{theorem}

\begin{proof}
Work in the Hurwitz integral ring of formal series
\[
 \mathcal H_p=\left\{F\in\Q_p[[t]]:
            d![t^d]F\in\Z_p\text{ for every }d\right\}.
\]
For $D=d/dt$, let $B_p=(D^{p^2}-D^p)/p$.  Call $F$ flat if
$B_p^sF\in\mathcal H_p$ for every $s\ge0$.
We will prove that $j!\mathcal P_{N,j}$ is flat at every address.  The main
difficulty is to include the addresses at which the usual unit-pivot
argument fails.

Set
\[
 H_c(t)=\sum_{n\ge0}\frac{t^{Nn+c}}{(Nn+c)!},\qquad
 A_c(t)=\sum_{n\ge0}\frac{t^{Nn}}{(Nn+c)!},\qquad
 R_c(t)=\frac{A_{c-1}(t)}{A_c(t)}\quad(1\le c<N).
\]
The denominator of this last quotient has nonzero constant coefficient.
The section identities give
\[
 R_cH_c=tH_{c-1},\qquad R_c(0)=c,\qquad
 j!\mathcal P_{N,j}=\mathcal P_{N,0}\prod_{c=1}^{j}R_c.
\]
Cancellation of consecutive source sections proves the product identity
and ties the ratios directly to the family under study.

At the singular addresses $c=ap$, $1\le a<m$, use the weighted ring
\[
 \mathcal W_p=\left\{\sum_{d\ge0}a_dT^d:
          p^{\lfloor d/p\rfloor}a_d\in\Z_p\right\}.
\]
The key estimate is $R_{ap}(\log(1+T))\in\mathcal W_p$.  Its arithmetic
input comes from the cyclic group ring $\Z[\Z/N\Z]$.
For its canonical element $x$ and $I=(p,x^{p-1}-1)$, the relation
$(x^{p-1})^p=1$ implies $I^p\subseteq(p)I$.  Induction yields
$I^q\subseteq(p^{v_p(q!)})$.  Grouping the falling factorial
$(x)_d=\prod_{i=0}^{d-1}(x-i)$ into consecutive blocks of length $p$
then gives an integral residue-coordinate factor
$p^{v_p(\lfloor d/p\rfloor!)}$.  The identity
\[
 v_p(d!)=\lfloor d/p\rfloor+v_p(\lfloor d/p\rfloor!)
\]
accounts for precisely the weight in $\mathcal W_p$ after dividing by $d!$.
Comparison of the section and group-ring recurrences identifies these
coordinates with the logarithmically substituted section.

To pass from the section to the singular ratio, factor out its initial zero:
for $c=ap$, the section has the form
\[
 H_c(\log(1+T))=\frac{T^c}{c!}\,V(T),\qquad V(0)=1,
 \qquad V\in\mathcal W_p.
\]
The weight of $c!$ compensates for the coefficient shift by $c$.
The weighted ring is closed under inversion of series with constant
coefficient one and under $L_T=\log(1+T)(1+T)d/dT$.  With
$\ell(T)=\log(1+T)/T$, interpreted by coefficient shifting, differentiation
of the displayed factorization gives
\[
 R_c(\log(1+T))=c(1+T)\ell(T)+(L_TV)V^{-1}.
\]
Since $p\ell(T)\in\mathcal W_p$ and $p\mid c$, both terms are weighted.
This is the singular estimate
\lean{KomorebiPaperIX.SingularRatioAmice.paper_singular_amice_estimate}.

To deduce flatness from the weighted estimate, use
$\Theta=(1+T)d/dT$.  Both $\Theta$ and
$(\Theta^{p^2}-\Theta^p)/p$ preserve $\mathcal W_p$.
For the first operator, its coefficient formula is
$(d+1)a_{d+1}+da_d$; whenever the floor weight rises, the factor $d+1$
provides the missing $p$.  For the second operator, the proof uses its
full Leibniz correction and the scalar congruence
$p^2\mid b^{p^2}-b^p$.  Weighted monomials factor as
\[
 a_dT^d=(p^{\lfloor d/p\rfloor}a_d)
          (T^p/p)^{\lfloor d/p\rfloor}T^{d\bmod p}.
\]
The operator preserves these monomials and hence finite weighted
polynomials.  Each fixed output coefficient depends on only finitely many
input coefficients, which extends the result to the formal series.
Substitution $T=e^t-1$ intertwines $D$ with $\Theta$ and maps
$\mathcal W_p$ into $\mathcal H_p$, because
\[
 d![t^d]F(e^t-1)=[T^0]\Theta^dF.
\]
It follows that the singular ratios are flat.  This transfer is proved in
\lean{KomorebiPaperIX.WeightedDworkTransfer.singular_ratio_flat}.

At a nonsingular address $2\le c<N$ with $p\nmid c$, differentiation of the section
identities gives
\[
 tR_c'=R_c(1+R_{c-1}-R_c).
\]
Writing $q=R_{c-1}-(c-1)$ and $y=R_c-c$, this becomes
\[
 (tD+c)y=cq+yq-y^2.
\]
On degrees divisible by $p$, the inverse of $tD+c$ divides the degree-$d$
coefficient by the unit $d+c$.  Starting with $y_0=0$, form
\[
 y_{s+1}=(tD+c)^{-1}(cq+y_sq-y_s^2).
\]
Flatness and support are preserved. Indeed, put $E_c=tD+c$ and $C_p=pD^{p^2}-D^p$.
The commutator is $[B_p,E_c]=C_p$, and iteration gives
\[
 B_p^rE_cy=E_cB_p^ry+rC_pB_p^{r-1}y.
\]
For the inverse equation $E_cy=f$, induction on $r$ makes the
correction term Hurwitz-integral. Both operators preserve support in
degrees divisible by $p$. On that support the degree multiplier $d+c$
is a unit, so integrality of $E_cB_p^ry$ implies integrality of
$B_p^ry$. Thus the inverse preserves full flatness, as well as support.
 Since the difference of two right
sides factors as $(z-y)(q-z-y)$ and its second factor has zero constant
term, agreement with the solution gains one coefficient at each iteration.
Consequently $[t^d]y_s=[t^d]y$ for $d\le s$.
Every coefficient of $B_p^r y$ can therefore be read from a sufficiently
long finite iterate.  The argument is entirely formal.
This is the content of
\lean{KomorebiPaperIX.RiccatiFlatPropagation.riccati_solution_flat}.

The standard window supplies the initial ratio $R_1$.  The weighted
estimate handles $p\mid c$, and Riccati propagation handles $p\nmid c$.
Flatness is closed under products, so the exact factorization establishes
flatness of $j!\mathcal P_{N,j}$.  Extracting its supported Hurwitz coefficients
and applying the proved saturation and step-multiple lemmas yields
\eqref{eq:arith-divisor-integral}--\eqref{eq:arith-divisor-bound}.

Here is the coefficient extraction explicitly. If $a_d=d![t^d]F$,
then
\[
 d![t^d]B_p^kF=p^{-k}\Delta_{p(p-1)}^ka_{d+pk}.
\]
For $F=j!\mathcal P_{N,j}$ and $k\le mn$, take $d=mpn-pk$.
Support in multiples of $mp$ turns this into
$p^{-k}\Delta_h^kU_{N,j}(n)$, so flatness gives divisibility by $p^k$.
For $k>mn$, expand $\Delta_h^{k-mn}$ applied to
$\Delta_h^{mn}U_{N,j}(n)$; every shifted row is at least $n$, so
each summand is divisible by $p^{mn}$. Finally
$\Delta_{uh}=((I+\Delta_h)^u-I)$, whose $k$th power for $u>0$
is an integral combination of powers of $\Delta_h$ of order at least $k$.
The case $u=0$ gives zero for $k>0$, and order zero is the original
integral coefficient. This proves the bound for all $u$ and $k$.
Finally $\Delta_{uh}^kU_{N,j}=j!\Delta_{uh}^kC_{N,j}$ gives
\eqref{eq:arith-divisor-normalized}, with the zero case retained.
The resulting statements are
\lean{KomorebiPaperIX.AllResidueEuler.all_residue_flat},
\lean{KomorebiPaperIX.AllResidueEuler.global_unnormalized_bound}, and
\lean{KomorebiPaperIX.AllResidueEuler.global_normalized_precision}.
\end{proof}

\begin{example}[Why the factorial loss is necessary]
\label{ex:arith-normalization-loss}
Take $p=3$, $m=2$, $N=6$, and $j=3$.  The recurrence gives
\[
 C_{6,3}(1)=-\frac1{84},\qquad
 C_{6,3}(2)=\frac{703}{5460},\qquad
 C_{6,3}(2)-C_{6,3}(1)=\frac{64}{455}.
\]
The last difference has $3$-adic valuation zero.  Here $h=1$, so at
$n=1$ the order-one bound without factorial loss would require valuation
at least one.  The corrected normalized bound is sharp, since
\[
 v_3\bigl(C_{6,3}(2)-C_{6,3}(1)\bigr)=0=1-v_3(3!).
\]
Multiplication by $3!$ gives
\[
 U_{6,3}(2)-U_{6,3}(1)=\frac{384}{455},\qquad
 v_3\bigl(U_{6,3}(2)-U_{6,3}(1)\bigr)=1,
\]
consistent with the unnormalized theorem.  Thus the factorial correction
is already necessary at a single step in the all-address parameter range.
These exact values and the refutation of a universal bound without loss
are formalized in
\lean{KomorebiPaperIX.Normalization.normalized_first},
\lean{KomorebiPaperIX.Normalization.normalized_second},
\lean{KomorebiPaperIX.Normalization.normalized_difference},
\lean{KomorebiPaperIX.Normalization.difference_valuation},
\lean{KomorebiPaperIX.Normalization.factorial_loss_is_sharp},
\lean{KomorebiPaperIX.Normalization.unnormalized_difference},
\lean{KomorebiPaperIX.Normalization.unnormalized_difference_valuation}, and
\lean{KomorebiPaperIX.Normalization.universal_no_loss_refuted}.
\end{example}

\begin{corollary}[A common period in the standard window]
\label{cor:arith-window-period}
Let $p$ be prime, $m>0$, $m\mid p-1$, and $0\le j<p$.  Put
$h=(p-1)/m$.  For every $n\ge1$ and $u\in\N$,
\[
 \overline C_m(n+uh)=\overline C_m(n),\qquad
 C_{mp,j}(n+uh)\equiv C_{mp,j}(n)\pmod p.
\]
Thus $h$ is a sufficient period from row one, common to all addresses
in the standard window.  It need not be the least period.
\end{corollary}

\begin{proof}
Since $j<p$, the factorial $j!$ is a unit in $\Z_p$.  Also $mn\ge1$.
The order-one case of Theorem~\ref{thm:arith-divisor-base} therefore gives
valuation at least one for the normalized difference at every step $uh$.
This proves the congruence, and Proposition~\ref{prop:arith-lucas-window}
identifies its reduction with the equality for $\overline C_m$.
The preserved Lean development also proves this period directly from the
Hurwitz inverse in characteristic $p$:
\lean{KomorebiPaperIX.WindowPeriod.uniform_standard_window_period},
\lean{KomorebiPaperIX.WindowPeriod.multiple_period}, and
\lean{KomorebiPaperIX.WindowPeriod.standard_window_period_divisibility}.
\end{proof}

\begin{lemma}[Prime-power period from finite differences]
\label{lem:arith-newton-period}
Let $R$ be a commutative ring, $p$ a prime, $f:\N\to R$, $h,n\in\N$,
and $r\ge1$.  If
\[
 p^{\min(k,r)}\mid\Delta_h^k f(n)\quad\text{for every }k\in\N,
\]
then $p^r\mid f(n+p^{r-1}h)-f(n)$.
\end{lemma}

\begin{proof}
The finite Newton expansion gives
\[
 f(n+p^{r-1}h)-f(n)
 =\sum_{k=1}^{p^{r-1}}\binom{p^{r-1}}k\Delta_h^k f(n).
\]
Terms with $k\ge r$ have the required factor from the difference.
For $1\le k<r$, the identity
\[
 k\binom{p^{r-1}}k=p^{r-1}\binom{p^{r-1}-1}{k-1}
\]
and $v_p(k)<k$ imply
$p^r\mid\binom{p^{r-1}}k p^k$ in $\Z$.
Mapping that integer divisibility to $R$ proves the claim term by term.
The argument includes $r=1$.
The declaration is
\lean{KomorebiPaperIX.PeriodFromDifferences.prime_power_period}.
\end{proof}

\begin{proof}[Proof of Theorem~\ref{thm:arith-nishibuchi}]
On the divisor branch, $m\mid q$ and the packet step $q/m$ is an integer
multiple of $h=(p-1)/m$.  If $r\le mn$, Theorem~\ref{thm:arith-divisor-base}
supplies the hypotheses of Lemma~\ref{lem:arith-newton-period}.
Its packet displacement is $p^{r-1}q/m$, which converts to the original
degree as follows:
\[
 mp\left(n+p^{r-1}\frac qm\right)=mpn+qp^r.
\]
The choice $n\ge r$ implies $mn\ge r$ because $m\ge1$.

For an odd prime, $2\mid p-1$, so the option $m=2$ introduces no further
case.  The sole remaining branch is $p=2,m=2,N=4$.  The all-row theorem
proved next gives $2^k\mid\Delta_1^kC_{4,j}(n)$ at every row and address.
After multiplying by $j!$, Lemma~\ref{lem:arith-newton-period} gives packet
period $2^{r-1}$ and degree displacement $4\cdot2^{r-1}=2^{r+1}=qp^r$.
It also gives integrality.  Hence the complete statement follows.
The formal degree conversion is
\lean{KomorebiPaperIX.NishibuchiCompletion.original_degree_shift};
the all-parameter degree theorem is
\lean{KomorebiPaperIX.NishibuchiCompletion.full_period_degree}.
\end{proof}

\subsection{All-row dyadic differences and displacement amplification}

For the remainder of the four-packet discussion abbreviate
$C_j=C_{4,j}$ and $U_j=U_{4,j}$.  Removing the original row cap requires
an unshifted coefficient operator.

\begin{theorem}[All-row dyadic higher differences]\label{thm:arith-dyadic-differences}
For $0\le j<4$ and all $u,k,n\in\N$,
\begin{align}
 C_j(n)&\in\Z_2,\\
 \Delta_u^kC_j(n)&\in 2^{k(1+\nu_2(u))}\Z_2,\label{eq:arith-allrow-c}\\
 \Delta_u^kU_j(n)&\in
 2^{k(1+\nu_2(u))+v_2(j!)}\Z_2.\label{eq:arith-allrow-u}
\end{align}
There is no restriction $k\le n$ and no minimum with the row $n$.
\end{theorem}

\begin{proof}
The Hurwitz coefficients of the normalized source $K_{4,j}$ at degrees
$4n$ are
\[
 g_j(n)=\binom{4n+j}{4n}^{-1}.
\]
Their explicit expressions are
\[
 \begin{aligned}
 g_0(n)&=1,&
 g_1(n)&=\frac1{4n+1},\\
 g_2(n)&=\frac1{(4n+1)(2n+1)},&
 g_3(n)&=\frac3{(4n+1)(2n+1)(4n+3)}.
 \end{aligned}
\]
Every denominator is odd.  An affine sequence with even slope satisfies
$2^k\mid\Delta_1^k f(n)$ for every $k,n$.  The exact product identity
\[
 \Delta_1^s(fg)(n)=\sum_{i=0}^s\binom si
          \Delta_1^if(n)\Delta_1^{s-i}g(n+i)
\]
proves closure of this property under products.  Applying the formula
with the factors interchanged to $fg=1$ isolates $f(n+s)\Delta_1^sg(n)$.
Multiplication by $f(n+s)^{-1}$ and induction then prove closure under
pointwise unit inverses.
Thus all four $g_j$ have the property.  The explicit formulas and difference
bounds are formalized in
\lean{KomorebiPaperIX.DyadicFourKernel.source_coe} and
\lean{KomorebiPaperIX.DyadicFourKernel.source_kummer}.

For an integer charge $c$, put
\[
 A_c=\frac{D^4-c}{2}.
\]
A series is flat of charge $c$ when applying any power $A_c^s$ to it gives
an element of $\mathcal H_2$.  The required product formula is
\begin{equation}\label{eq:arith-charged-product}
 A_{c+d}(FG)=(A_cF)G+F(A_dG)
       +2F'G'''+3F''G''+2F'''G'.
\end{equation}
All middle coefficients are integral.  Derivatives preserve
$\mathcal H_2$ and commute with $A_c$.  Induction on the number of
iterations in \eqref{eq:arith-charged-product} proves that multiplication
adds charges.  If $FG=1$, the same formula yields
\[
 A_{-c}G=-G^2A_cF-G(2F'G'''+3F''G''+2F'''G').
\]
Induction, starting from the independently established integrality of
the reciprocal $G$, proves that inversion negates the charge.
Finally $A_{c+2}=A_c-1$, so increasing a charge by two preserves flatness.
These are
\lean{KomorebiPaperIX.ChargedDworkClosure.B_mul},
\lean{KomorebiPaperIX.ChargedDworkClosure.flat_inverse}, and
\lean{KomorebiPaperIX.ChargedDworkClosure.flat_charge_shift}.

Packing $g_j$ into degrees divisible by four gives $K_{4,j}$.
Its source property is precisely charge-one flatness.  Inversion gives
charge $-1$, and the charge shift restores charge one for the
reciprocal $\mathcal P_{4,j}$.  If $\operatorname{dec}(F)(d)=d![t^d]F$, the exact
coefficient identity is
\[
 \operatorname{dec}(A_1^sF)(d)
       =2^{-s}\Delta_4^s\operatorname{dec}(F)(d).
\]
At $d=4n$, packing turns the right side into
$2^{-s}\Delta_1^sC_j(n)$.  This extraction requires no negative degree
and therefore holds at every $n$, including zero.  It proves the unit-step
bound \lean{KomorebiPaperIX.DyadicFourAllRows.base_step_bound}.

To incorporate $v_2(u)$, expand a displacement using finite Newton
coefficients.  For any prime $p$, $u>0$, and $1\le a\le u$,
\[
 p^{v_p(u)+1}\mid\binom ua p^a,
\]
as follows from $a\binom ua=u\binom{u-1}{a-1}$ and $v_p(a)<a$.
Repeated displacement, retaining the remaining base-difference order
at each induction step, gives a factor $p^{k(1+v_p(u))}$ from $k$
displacements.  This induction is formalized in
\lean{KomorebiPaperIX.StepAmplification.repeated_precision}.
Applying it with $p=2$ proves \eqref{eq:arith-allrow-c}; multiplication
by $j!$ proves \eqref{eq:arith-allrow-u}.  Zero step and zero order follow
from their definitions.  The final statements are
\lean{KomorebiPaperIX.DyadicFourFinalStrengthening.normalized_step_bound}
and
\lean{KomorebiPaperIX.DyadicFourFinalStrengthening.unnormalized_step_bound}.
\end{proof}

\begin{corollary}[Step-sensitive divisor-branch bound]\label{cor:arith-amplified}
Under the hypotheses of Theorem~\ref{thm:arith-divisor-base}, for all
$u,k,n\in\N$,
\[
 \Delta_{u(p-1)/m}^{k}U_{mp,j}(n)
       \in p^{\min(k(1+\nu_p(u)),mn)}\Z_p.
\]
For $C_{mp,j}$ the corresponding valuation lower bound is
$\min(k(1+\nu_p(u)),mn)-v_p(j!)$.
\end{corollary}

\begin{proof}
Apply the same repeated-displacement theorem with the fixed cap $mn$ to
\eqref{eq:arith-divisor-bound}.  The argument preserves this cap and
improves the contribution from the order of the difference.
Division by the fixed factorial gives the rational valuation statement,
with the zero case included.
These are
\lean{KomorebiPaperIX.AllResidueStepAmplification.unnormalized_precision}
and
\lean{KomorebiPaperIX.AllResidueStepAmplification.normalized_precision}.
\end{proof}

\begin{corollary}[Dyadic periods from row zero]\label{cor:arith-allrow-period}
For $0\le j<4$, $r\ge1$, and $n\ge0$,
\[
 U_j(n+2^{r-1})\equiv U_j(n)\pmod{2^r}.
\]
For $j=2,3$ and $r\ge2$, $2^{r-2}$ is also a sufficient packet period.
The respective original degree displacements are $2^{r+1}$ and $2^r$.
\end{corollary}

\begin{proof}
Use Lemma~\ref{lem:arith-newton-period} and the unit-step case of
Theorem~\ref{thm:arith-dyadic-differences}.  At $j=2,3$ the factorial
contributes an additional power of two.  The original degree versions are
\lean{KomorebiPaperIX.DyadicFourFinalStrengthening.all_rows_period_degree}
and
\lean{KomorebiPaperIX.DyadicFourFinalStrengthening.all_rows_half_period_degree}.
The next theorem sharpens these sufficient periods to a complete
classification.
\end{proof}

\subsection{Exact displacement valuations and least packet periods}

\begin{lemma}[Finite Newton dominance]\label{lem:arith-dominance}
Let $f:\N\to\Q$ and write $a_k=\Delta_1^kf(0)$.  Suppose that $c\in\N$,
$a_1\ne0$, $v_2(a_1)=c$, and
\[
 v_2(a_k)\ge c+1+\lfloor\log_2k\rfloor\qquad(k\ge2).
\]
Then for every $n\ge0$ and $T>0$,
\[
 f(n+T)-f(n)\ne0,\qquad
 v_2\bigl(f(n+T)-f(n)\bigr)=v_2(T)+c.
\]
\end{lemma}

\begin{proof}
Finite Newton reconstruction at an arbitrary row gives
\[
 \Delta_1^kf(n)=\sum_{i=0}^{n}\binom ni a_{k+i}.
\]
The monotonicity of $\lfloor\log_2k\rfloor$ therefore preserves the
higher-coefficient bound after translation.  In addition,
$\Delta_1f(n)-a_1\in2^{c+1}\Z_2$, so the first difference is nonzero
of exact valuation $c$.

Now expand the displacement:
\[
 f(n+T)-f(n)=T\Delta_1f(n)
          +\sum_{k=2}^{T}\binom Tk\Delta_1^kf(n).
\]
For $1\le k\le T$, the binomial denominator identity implies
\[
 v_2\binom Tk\ge v_2(T)-v_2(k)
            \ge v_2(T)-\lfloor\log_2k\rfloor.
\]
Every term of the remaining sum consequently has valuation at least
$v_2(T)+c+1$.  Its sum cannot cancel the leading term, which has exact
valuation $v_2(T)+c$.  This proves nonvanishing and the claimed valuation
using only finite sums.
The general theorem is
\lean{KomorebiPaperIX.NewtonDominance.exact_displacement}.
\end{proof}

\begin{theorem}[Exact four-packet displacement law]\label{thm:arith-exact}
Set
\[
 (c_0,c_1,c_2,c_3)=(1,1,4,2),\qquad
 (d_0,d_1,d_2,d_3)=(1,1,5,3).
\]
For $0\le j<4$, $n\ge0$, and every positive integer $T$,
\begin{align*}
 C_j(n+T)-C_j(n)&\ne0,&
 v_2\bigl(C_j(n+T)-C_j(n)\bigr)&=v_2(T)+c_j,\\
 U_j(n+T)-U_j(n)&\ne0,&
 v_2\bigl(U_j(n+T)-U_j(n)\bigr)&=v_2(T)+d_j.
\end{align*}
\end{theorem}

\begin{proof}
The first Newton coefficients, derived from
\eqref{eq:arith-triangular}, are
\[
 a_{0,1}=-2,\quad a_{1,1}=-\frac65,\quad
 a_{2,1}=-\frac{16}{15},\quad a_{3,1}=-\frac{36}{35}.
\]
Their exact valuations are $c_j$.  Theorem~\ref{thm:arith-dyadic-differences}
gives $v_2(a_{j,k})\ge k$ for all orders.  To obtain the stronger
logarithmic envelope required by Lemma~\ref{lem:arith-dominance}, only
a finite initial segment needs improvement.  The recurrence proves the
following exact valuations:
\begin{center}
\begin{tabular}{cll}\toprule
Address $j$ & Further orders $k$ & $v_2(a_{j,k})$ in that order\\\midrule
$0$ & $2$ & $3$\\
$1$ & $2$ & $3$\\
$2$ & $2,3,4,5,6$ & $6,9,13,16,22$\\
$3$ & $2,3,4$ & $8,8,15$\\\bottomrule
\end{tabular}
\end{center}
For all subsequent $k$, the elementary inequality
$c_j+1+\lfloor\log_2k\rfloor\le k$ completes the envelope.
The entries follow by exact rational computation from the defining
recurrence, and their valuations are proved in Lean.  Together with the
inequality for all subsequent orders, they establish the bound at every
order.
The uniform envelope is
\lean{KomorebiPaperIX.DyadicFourPrefix.tail_precision}.

Apply Lemma~\ref{lem:arith-dominance} to $C_j$.  Multiplication by the
nonzero factorial $j!$ adds $v_2(j!)$ and preserves nonzeroness.
Hence $d_j=c_j+v_2(j!)$, giving the four valuation gains for the original
coefficients.
The declarations are
\lean{KomorebiPaperIX.DyadicFourLeastPeriods.normalized_exact_displacement}
and \lean{KomorebiPaperIX.DyadicFourLeastPeriods.exact_displacement}.
\end{proof}

\begin{corollary}[All positive periods and their least element]
\label{cor:arith-least-period}
For $r\in\N$ and $0\le j<4$, a positive integer $T$ is a packet period
of $U_j$ modulo $2^r$ if and only if
\[
 2^{\max(0,r-d_j)}\mid T.
\]
Here a packet period means $U_j(n+T)\equiv U_j(n)\pmod{2^r}$ for every
$n\ge0$.  Consequently the least positive packet periods are
\[
 \begin{array}{c|cccc}
 j&0&1&2&3\\\hline
 \text{least period}&2^{\max(0,r-1)}&2^{\max(0,r-1)}&
                 2^{\max(0,r-5)}&2^{\max(0,r-3)}.
 \end{array}
\]
\end{corollary}

\begin{proof}
By the exact displacement law, the congruence at any row is equivalent to
$v_2(T)+d_j\ge r$.  This is exactly the displayed integer
divisibility criterion.  Its least positive solution is the indicated
power of two, which also covers $r=0$ and $r\le d_j$.
The all-period and least-element statements are
\lean{KomorebiPaperIX.DyadicFourLeastPeriods.period_iff_dvd} and
\lean{KomorebiPaperIX.DyadicFourLeastPeriods.least_packet_period}.
\end{proof}

For example, $U_2(n)\equiv2\pmod{32}$ and $U_3(n)\equiv6\pmod8$
for every row.  At precision $2^{10}$ the least packet periods are
$512,512,32,128$.  A packet displacement $T$ has original degree
displacement $4T$, as proved by
\lean{KomorebiPaperIX.DyadicFourLeastPeriods.original_degree_exact_displacement}.
Minimality here concerns packet shifts.  Shifts of the full degree-indexed
sequence, including its zero coefficients off the packet support, require
a separate analysis.  Exact valuations of arbitrary higher-order differences
are also a separate question from the first-displacement law.

\subsection{Classical odd-length coordinates and their congruence ladder}

The classical signed Euler numbers have generating function
$\operatorname{sech}t=\sum_{n\ge0}E_{2n}t^{2n}/(2n)!$.
They are the $N=2$, $j=0$ member of
\eqref{eq:arith-euler-egf}.  The coordinate theorem reconstructs
them from finite alternating sums of odd powers.

\begin{definition}[Odd-length source and inverse coefficients]
\label{def:arith-odd-source}
For a positive odd integer $p$, define
\[
 S_d(p)=\sum_{a=0}^{p-1}(-1)^a(2a+1)^d.
\]
Let $\gamma_m$ be the integral Hurwitz coefficients of $\operatorname{sech}^2$:
\[
 \operatorname{sech}^2t=\sum_{m\ge0}\gamma_m\frac{t^{2m}}{(2m)!},
 \qquad
 \gamma_m=\sum_{b=0}^{2m}\binom{2m}{b}E_bE_{2m-b},
\]
where $E_b=0$ for odd $b$.  Then $\gamma_0=1$.  Define
\[
 T^{\mathrm{odd}}_{p,M}(n)=\sum_{m=0}^{\min(M,n)}
       \binom{2n}{2m}\gamma_m p^{2m}S_{2(n-m)}(p).
\]
Here $p$ need only be odd; primality is needed only to interpret the
result as a prime-adic congruence ladder.
\end{definition}

The source and inverse coefficients are
\lean{KomorebiPaperIX.OddPower.source} and
\lean{KomorebiPaperIX.EulerCoefficients.inverseEven}; the truncation is
\lean{KomorebiPaperIX.OddEuler.truncation}.

The same inverse coefficients can be written in the signed zigzag notation
of the original manuscript.  Define $Z_r$ by
\[
 \sec t+\tan t=\sum_{r\ge0}Z_r\frac{t^r}{r!}.
\]
Then
\begin{equation}\label{eq:arith-signed-zigzag}
 \gamma_m=(-1)^m Z_{2m+1}\qquad(m\in\N).
\end{equation}
Indeed, the odd coefficients come from $\tan t$, whose derivative is
$\sec^2t$.  Changing the sign of the squared variable in the even series
for $\operatorname{sech}^2t$ gives $\sec^2t$, so its degree-$2m$
Hurwitz coefficient is $(-1)^m\gamma_m$.  Differentiation identifies
this coefficient with $Z_{2m+1}$, proving the formula.  The Lean bridge is
\lean{KomorebiPaperIX.Zigzag.inverse_is_signed_zigzag}; it uses the
formal generating-function definition of the zigzag numbers.

\begin{theorem}[Classical odd-length reconstruction]
\label{thm:arith-odd-reconstruction}
For every positive odd integer $p$ and every $n\in\N$,
\[
 E_{2n}=\sum_{m=0}^{n}\binom{2n}{2m}\gamma_m p^{2m}S_{2(n-m)}(p).
\]
Equivalently, $E_{2n}=\sum_{q=0}^n V_{nq}S_{2q}(p)$ for the lower
triangular matrix
\[
 V_{nq}=\binom{2n}{2q}\gamma_{n-q}p^{2(n-q)}\quad(q\le n),
 \qquad V_{nq}=0\quad(q>n).
\]
It has $V_{nn}=1$ and $p^{2(n-q)}\mid V_{nq}$ when $q\le n$.
\end{theorem}

\begin{proof}
The exponential generating series of the source is the finite sum
\[
 \mathcal S_p(t)=\sum_{a=0}^{p-1}(-1)^a e^{(2a+1)t}.
\]
Multiplication by $\cosh t$ telescopes.  Since $p$ is odd,
\[
 \mathcal S_p(t)=\frac{1+e^{2pt}}{2}\operatorname{sech}t.
\]
Taking its even part gives
\[
 \sum_{n\ge0}S_{2n}(p)\frac{t^{2n}}{(2n)!}
       =\cosh^2(pt)\operatorname{sech}t.
\]
Multiplication by the reciprocal $\operatorname{sech}^2(pt)$
yields $\operatorname{sech}t$.  The EGF product rule now gives the
reconstruction coefficient by coefficient.  Replacing $m$ by $n-q$
produces the displayed lower triangular row.  Its diagonal and weights
follow from $\gamma_0=1$ and the explicit factor $p^{2(n-q)}$.
The source-series identity is
\lean{KomorebiPaperIX.OddPower.odd_source_even_egf}; the reconstruction
and row statement are
\lean{KomorebiPaperIX.OddEuler.odd_euler_reconstruction} and
\lean{KomorebiPaperIX.OddEuler.row_action}.
\end{proof}

\begin{corollary}[Odd-prime truncation ladder]\label{cor:arith-odd-ladder}
For every positive odd integer $p$ and all $M,n\in\N$,
\[
 p^{2M+2}\mid E_{2n}-T^{\mathrm{odd}}_{p,M}(n)\quad\text{in }\Z.
\]
If $M\ge n$, the truncation equals $E_{2n}$ exactly.  In particular,
for an odd prime $p$ this is a congruence modulo $p^{2M+2}$ in $\Z_p$.
\end{corollary}

\begin{proof}
Every omitted summand has $m\ge M+1$, hence contains the integer factor
$p^{2m}$ divisible by $p^{2M+2}$.  All its other factors are integral.
The finite sum of these omitted terms proves the assertion, with the
empty-tail case giving exact equality.  These are
\lean{KomorebiPaperIX.OddEuler.odd_euler_supercongruence} and
\lean{KomorebiPaperIX.OddEuler.full_depth_exact}.
\end{proof}

The reconstruction and its truncation bound already hold over the integers.
For the dyadic construction we use a different source.  Its original bound
of $2^{2M+2}$ follows from integral reciprocal coefficients, and the next
subsection sharpens that bound to the optimal uniform exponent $4M+3$.

\subsection{The optimal classical dyadic edge tower}

Return now to the classical signed Euler numbers
\[
 \operatorname{sech}t=\frac1{\cosh t}
       =\sum_{n\ge0}E_{2n}\frac{t^{2n}}{(2n)!}.
\]
Thus $E_{2n}=E_{2n}^{(2,0)}$ in the convention
\eqref{eq:arith-euler-egf}.  In particular $E_0=1$, $E_2=-1$, and $E_4=5$.
The following construction concerns these $N=2$ coefficients, not $C_{4,0}$.

\begin{definition}[Dyadic truncation]\label{def:arith-truncation}
Define $\beta_m$ by the formal unit-series identity
\[
 \sum_{m\ge0}\beta_m z^{2m}=\left(\frac{z}{\sinh z}\right)^2,
 \qquad \beta_0=1,
\]
where the quotient means the inverse of the unit series
$(\sinh z/z)^2$.  Define
\[
 A(n)=\frac{3^{2n+2}-1}{4(2n+1)(2n+2)},\qquad
 S_m(n)=\beta_m2^{2m}\frac{(2n)!}{(2n-2m)!}A(n-m)
       \quad(0\le m\le n),
\]
and
\begin{equation}\label{eq:arith-tm}
 T_M(n)=\sum_{m=0}^{\min(M,n)}S_m(n)\qquad(M,n\in\N).
\end{equation}
\end{definition}

These definitions correspond to
\lean{KomorebiPaperIX.Dyadic.beta},
\lean{KomorebiPaperIX.Dyadic.sourceEven}, and
\lean{KomorebiPaperIX.DyadicReconstruction.truncation}.
The formal summand uses a falling factorial, and
\lean{KomorebiPaperIX.DyadicReconstruction.falling_factorial_fraction}
proves its equality with the quotient displayed above.

\begin{lemma}[Exact reconstruction and source integrality]
\label{lem:arith-reconstruction}
For every $n\in\N$,
\[
 E_{2n}=\sum_{m=0}^{n}S_m(n).
\]
Every $\beta_m$ is $2$-adically integral and every $A(n)$ is a nonzero
$2$-adic unit.  If $M\ge n$, then $T_M(n)=E_{2n}$.
\end{lemma}

\begin{proof}
The ordinary kernel coefficients of $(\sinh z/z)^2$, in $z^2$, are
\[
 \frac{2^{2m+1}}{(2m+2)!}.
\]
Legendre's strict factorial bound
$v_2((2m+2)!)<2m+2$ makes them integral.  Their constant coefficient is
one.  The coefficient recurrence for its reciprocal therefore
proves integrality of every $\beta_m$.
For the source factor, lifting the exponent at the even power $2n+2$
gives
\[
 v_2(3^{2n+2}-1)=2+v_2(2n+2).
\]
Its denominator has the same valuation, since $2n+1$ is odd.  Thus
$A(n)$ is a nonzero unit.

Its exponential generating function is
\[
 \mathcal A(t)=\sum_{n\ge0}A(n)\frac{t^{2n}}{(2n)!}
             =\frac{\cosh(3t)-\cosh t}{4t^2}.
\]
The quotients by powers of $t$ are well defined because the required
initial coefficients vanish.  The hyperbolic identity
\[
 \mathcal A(t)\cosh t=\left(\frac{\sinh(2t)}{2t}\right)^2
\]
then gives the exact formal-series product
\[
 \left(\sum_{m\ge0}\beta_m(2t)^{2m}\right)\mathcal A(t)
       =\operatorname{sech}t.
\]
Taking the factorial-scaled coefficient of $t^{2n}$ produces the stated
finite sum.  Full depth leaves no omitted terms.
The relevant declarations are
\lean{KomorebiPaperIX.Dyadic.beta_integral},
\lean{KomorebiPaperIX.Dyadic.source_valuation},
\lean{KomorebiPaperIX.DyadicSeries.dyadic_reconstruction_series}, and
\lean{KomorebiPaperIX.DyadicReconstruction.dyadic_reconstruction}.
\end{proof}

\begin{lemma}[Exact weighted inverse coefficient]\label{lem:arith-kappa}
Let $\kappa_m=\beta_m2^{2m}(2m)!$.  For every $m\ge1$,
\[
 \kappa_m=-2^{4m}(2m-1)B_{2m}\ne0,\qquad
 v_2(\kappa_m)=4m-1,
\]
where $B_{2m}$ is the usual even Bernoulli number.
\end{lemma}

\begin{proof}
Let $Q(z)=z\coth z$, the regular formal series with constant coefficient
one.  The elementary differential identity
\[
 \left(\frac z{\sinh z}\right)^2=Q(z)-zQ'(z)
\]
and the Bernoulli expansion
\[
 Q(z)=1+\sum_{m\ge1}\frac{2^{2m}B_{2m}}{(2m)!}z^{2m}
\]
give
$\beta_m=-2^{2m}(2m-1)B_{2m}/(2m)!$.
Multiplication by $2^{2m}(2m)!$ proves the identity for $\kappa_m$.
The Bernoulli denominator theorem gives $v_2(B_{2m})=-1$ for $m>0$: its denominator contains each prime $q$ with $(q-1)\mid 2m$ exactly once, including $2$\cite{DLMF24}.
Since $2m-1$ is odd, the valuation is exactly $4m-1$.
This also proves nonvanishing.  In Lean, the hyperbolic identity is proved
after clearing removable zeros.  The library Bernoulli theorem then
determines the valuation of $\kappa_m$.
The declarations are
\lean{KomorebiPaperIX.DyadicSharpTower.kappa_bernoulli} and
\lean{KomorebiPaperIX.DyadicSharpTower.kappa_valuation}.
\end{proof}

\begin{theorem}[Optimal uniform dyadic tower]\label{thm:arith-sharp-tower}
For all $M,n\in\N$,
\begin{equation}\label{eq:arith-sharp-tail}
 E_{2n}-T_M(n)\in2^{4M+3}\Z_2.
\end{equation}
At the first omitted row $n=M+1$, the error is nonzero and
\[
 v_2\bigl(E_{2M+2}-T_M(M+1)\bigr)=4M+3.
\]
Thus the exponent $4M+3$ cannot be increased uniformly in the Euler row.
\end{theorem}

\begin{proof}
Rewrite the summand exactly as
\[
 S_m(n)=\kappa_m\binom{2n}{2m}A(n-m).
\]
For $m>M$, its first factor is divisible by $2^{4m-1}$, the binomial
coefficient is integral, and $A(n-m)$ is a unit.  Hence every omitted
summand is divisible by $2^{4M+3}$.  Summing the finite tail proves
\eqref{eq:arith-sharp-tail}.  If $M\ge n$, the tail is empty and the
statement includes the resulting zero error.

At $n=M+1$ there is exactly one omitted summand.  Here
$\binom{2n}{2n}=1$ and $A(0)=1$, so the error is $\kappa_{M+1}$.
Lemma~\ref{lem:arith-kappa} gives its exact valuation $4M+3$ and
nonvanishing.  The statements are
\lean{KomorebiPaperIX.DyadicSharpTower.dyadic_supercongruence_sharp},
\lean{KomorebiPaperIX.DyadicSharpTower.first_omitted_error}, and
\lean{KomorebiPaperIX.DyadicSharpTower.first_omitted_error_valuation}.
\end{proof}

\begin{proposition}[Row-dependent summand weights]\label{prop:arith-row-weight}
For $1\le m\le n$, the summand is nonzero and
\[
 v_2(S_m(n))=4m-1+v_2\binom{2n}{2m}.
\]
More generally, if $L\in\N$ satisfies
\[
 L\le4m-1+v_2\binom{2n}{2m}\qquad\text{for every }M<m\le n,
\]
then $E_{2n}-T_M(n)\in2^L\Z_2$.
\end{proposition}

\begin{proof}
Use the same factorization of $S_m(n)$, the nonzero binomial coefficient
for $m\le n$, the exact valuation of $\kappa_m$, and the unit property
of $A(n-m)$.  The tail conclusion follows by addition in $2^L\Z_2$.
These are \lean{KomorebiPaperIX.DyadicSharpTower.summand_valuation} and
\lean{KomorebiPaperIX.DyadicSharpTower.dyadic_supercongruence_row_weight}.
\end{proof}

The minimum summand weight bounds the tail valuation from below, but
cancellation may raise that valuation.  An exact formula for the entire
tail at every row therefore remains open.  Uniform optimality follows
from the equality already established at $n=M+1$.

\subsection{Isometric continuation of the four dyadic packets}

The exact displacement law also determines the images of the packets
after continuous extension to the dyadic integers.

\begin{theorem}[Isometric packet bijections]\label{thm:arith-isometry}
For $0\le j<4$, define on natural numbers
\[
 f_j(n)=\frac{U_j(n)-j!}{2^{d_j}},\qquad (d_0,d_1,d_2,d_3)=(1,1,5,3).
\]
There is a unique continuous map $\Phi_j:\Z_2\to\Z_2$ whose value at
every $n\in\N$ equals the embedded rational $f_j(n)$.
It is an isometric bijection and $\Phi_j(0)=0$:
\[
 |\Phi_j(x)-\Phi_j(y)|_2=|x-y|_2\qquad(x,y\in\Z_2).
\]
In particular, every element of $\Z_2$ has exactly one preimage.
\end{theorem}

\begin{proof}
For $n>0$, the displacement law at row zero gives
$v_2(U_j(n)-j!)=v_2(n)+d_j$, so $f_j(n)\in\Z_2$.
The value at zero is zero.  Applying the same law to the smaller of two
unequal natural numbers shows
\[
 |f_j(n)-f_j(m)|_2=|n-m|_2.
\]
The natural numbers are dense in $\Z_2$, and the target is complete.
The map therefore extends continuously.  The distance identity extends
on the dense product of natural-number points, making the extension
an isometry.  Density also proves uniqueness among all continuous maps
with these prescribed values.

To prove surjectivity, use compactness.  For an isometry $g$ of a compact
metric space into itself, the sequence of iterates $g^n(x)$ has a
convergent subsequence.  Two sufficiently close iterates, with exponents
$a<b$, yield $d(x,g^{b-a}(x))$ equally small after removing $a$ isometric
iterations.  Since $b-a>0$, the latter point is in the image of $g$.
Thus the image is dense.  It is also compact and hence closed, so it is
the whole space.  Apply this to $\Phi_j$ on $\Z_2$.

The general ingredients are
\lean{KomorebiPaperIX.DensePadicIsometry.existsUnique_extensionEquiv} and
\lean{KomorebiPaperIX.CompactSelfIsometry.surjective}.
The packet specialization is
\lean{KomorebiPaperIX.DyadicFourIsometry.existsUnique_isometry};
uniqueness among continuous maps is
\lean{KomorebiPaperIX.DyadicFourIsometry.unique_continuous}.
\end{proof}

\begin{corollary}[Exact image of the original coefficients]
\label{cor:arith-affine-image}
The continuous extension of the original sequence is
\[
 \widetilde U_j(x)=j!+2^{d_j}\Phi_j(x).
\]
It is injective, its image is exactly $j!+2^{d_j}\Z_2$, and
\[
 |\widetilde U_j(x)-\widetilde U_j(y)|_2
     =2^{-d_j}|x-y|_2.
\]
Every point of that affine image has exactly one preimage in $\Z_2$.
\end{corollary}

\begin{proof}
The formula agrees with $U_j$ on every natural number.  Continuity and
the distance equality follow from the corresponding properties of
$\Phi_j$.  Its surjectivity gives precisely the displayed affine image,
and multiplication by the nonzero scalar $2^{d_j}$ preserves injectivity.
The declarations are
\lean{KomorebiPaperIX.DyadicFourIsometry.originalExtension_nat},
\lean{KomorebiPaperIX.DyadicFourIsometry.originalExtension_dist}, and
\lean{KomorebiPaperIX.DyadicFourIsometry.originalExtension_range}.
\end{proof}

The four rational packet sequences thus extend to bijections onto their
exact dyadic images.  The proof uses density, completeness, and the
displacement law.  Analyticity of these packet extensions would require
a further argument.  The analytic polynomial solution maps constructed
later rest on separate hypotheses and a separate proof.
